\section{引言}
\label{sec:intro}

部分子分布函数（PDF）描述强子内部部分子的动力学~\cite{Feynman:102074}。它们具有非微扰的本质，因此无法在微扰论中计算。
格点 QCD 是在大规模模拟中从第一性原理研究 PDF 的理想框架。然而，PDF 通常定义在光锥上，
这给标准的欧氏表述带来了困难。于是，格点 QCD 中的强子结构计算与其他在欧氏时空中可及的量联系在一起。
这导致了 PDF 的 Mellin 矩与核子形状因子的长期研究历史（近期综述见
\cite{Constantinou:2014tga, Constantinou:2015agp,Alexandrou:2015yqa, Alexandrou:2015xts, Syritsyn:2014saa}）。
在实践中，高阶矩的信噪比很小，使 PDF 的重构受到严重限制。此外，还存在与更低量纲算符的
幂发散混合这一不可避免的问题。因此，从矩重构 PDF 的任务实际上是不可行的。

\medskip
PDF 的 {\textit{从头计算}}（{\textit{ab initio}}）极为重要，因为它将是对 QCD 非微扰性质的严格检验。
缺少从第一性原理出发的 PDF 计算，是一个阻碍我们更深入理解核子结构的紧迫问题。
我们目前关于 PDF 的知识依赖于利用微扰论对实验数据所做的唯象拟合。这些参数化的 PDF 例如作为
正在运行的对撞机（尤其是 LHC）截面计算的输入，也用于规划未来的对撞机实验——这些实验本身就是对标准模型的检验。
然而，参数化并非没有歧义~\cite{Jimenez-Delgado:2013sma}。此外，还存在实验上无法到达的运动学区域。
大 Bjorken $x$ 区域便是其中之一，其巨大的不确定性在下夸克分布中最为显著
\cite{Accardi:2016qay,Alekhin:2017kpj}。横向性（transversity）PDF 则是另一个仅被唯象粗略约束的 PDF 例子。

\medskip
X.~Ji 提出了一种直接计算的先驱性方法~\cite{Ji:2013dva}。在该方法中，人们不再计算定义在光锥上的矩阵元，
而是计算包含有限长度 Wilson 线的费米子算符的矩阵元，其傅里叶变换定义了所谓的准部分子分布（quasi-PDF）。
这通过把 Wilson 线取在纯空间方向（约定为 $z$ 方向）、而非光锥上的 $+$ 方向来实现。
就强子运动学而言，核子动量通常沿该空间方向取。在较大但有限的动量下，准 PDF 可通过匹配手续
与光前 PDF 相联系~\cite{Xiong:2013bka,Chen:2016fxx}。Ji 的方法已经在文献
Refs.~\cite{Lin:2014zya,Gamberg:2014zwa,Alexandrou:2015rja,Chen:2016utp,Alexandrou:2016jqi} 中得到检验，
所有这些结果都很有前景，即它们在匹配手续之后给出正确的 PDF 形状。准 PDF 的某些性质，如核子质量依赖和
靶质量效应，也已通过其与横动量依赖分布函数（TMD）的关系得到分析~\cite{Radyushkin:2017ffo,Radyushkin:2016hsy}。
还应注意相关方法——赝部分子分布（pseudo-PDF）——的出现，它是光锥 PDF 向有限核子动量的另一种推广
\cite{Radyushkin:2017cyf,Orginos:2017kos}。文献 Refs.~\cite{Carlson:2017gpk,Briceno:2017cpo} 讨论了欧氏符号在准 PDF
计算中的作用（与光前 PDF 相比较），其中后者证明了由欧氏格点 QCD 得到的矩阵元与在闵氏空间中使用 LSZ 约化公式得到的相同。
尽管如此，准 PDF 的矩阵元含有发散，必须通过重正化予以消除，才能得到可与物理 PDF 比较的有意义结果。

\medskip
迄今为止，所有关于准 PDF 的工作都只考虑了裸矩阵元，因为重正化过程高度非平凡，直到最近才被处理。
特别地，与局域算符相比，Wilson 线算符的重正化出现了新的复杂性。其一，除对数发散之外，
还存在相对格点调节子 $a$ 的线性发散~\cite{Dotsenko:1979wR}，它在被消除之前使人无法取连续极限。
在单圈微扰论层面，该发散表现为线性发散，文献~\cite{Constantinou:2017sej} 对多种费米子和胶子作用量做了计算。
然而，非微扰地提取这一幂发散至关重要，这正是本工作的目标之一。这类算符带来新复杂性的另一特征是：
Dirac 结构的某些选择会表现出混合~\cite{Constantinou:2017sej}。

\medskip
证明这些矩阵元可以被重正化、特别是与 Wilson 线相联系的线性发散可以被消除，具有头等重要性。
没有重正化，整个准 PDF 方案就是不完整的，无法向理论和实验界提供任何有用信息。
文献 Refs.~\cite{Ma:2014jla,Ishikawa:2016znu} 提出了通过静态势消除线性发散的一些建议。
文献~\cite{Chen:2016fxx} 做了线性发散的单圈计算，并由此推动了无幂发散的改进准 PDF 的定义。
尽管如此，仍必须证明这样的步骤能够非微扰地完成。利用准 PDF 的核子矩阵元来提取线性发散系数的
另一种方法也在文献~\cite{Constantinou:2017sej} 中给出。文献~\cite{Monahan:2016bvm} 讨论了
利用梯度流压制线性发散的一种替代技术。最近有两篇文章讨论了对准 PDF 重正化采用辅助场形式
~\cite{Ji:2017oey,Green:2017xeu}，其中 Wilson 线被所引入辅助场的格林函数取代。
准 PDF 在微扰论各阶的可重正化性在文献~\cite{Ishikawa:2017faj} 中得到研究。

\medskip
本文首次提出了一种完全非微扰方式的准 PDF 具体重正化方法\footnote{在我们的工作提交之后，
所提出的重正化方案也被文献~\cite{Chen:2017mzz} 应用。}。
我们给出了该方法的具体处方，并展示了重正化矩阵元的例子。我们还讨论了如何消除非极化准 PDF
与 twist-3 标量算符之间的混合。我们采用 RI$'$ 重正化方案~\cite{Martinelli:1994ty}，
并利用文献~\cite{Constantinou:2017sej} 计算的单圈转换因子，把结果转换到参考标度 $\bar\mu{=}2$ GeV 的
$\MSb$ 方案。作为检验情形，由于螺旋度准 PDF 无混合，我们聚焦于它来展示结果。

\medskip
本文结构如下：第 2 节给出与核子矩阵元相关的理论设定，包括有混合和无混合情形下的重正化处方，
以及微扰计算的转换因子数据。第 3 节包括重正化函数的结果、$Z$ 因子系统不确定性的讨论、
螺旋度情形的重正化矩阵元，以及与光前 PDF 匹配的结果。最后我们进行总结并给出未来方向。
